\documentclass{article} 

\usepackage{graphicx}
\usepackage{subcaption}
\usepackage{float}
\usepackage{caption}

\begin{document}
\title{EAGLE BIG1 and BIG12 comparison}
\date{\today}
\maketitle

\section{Introduction}
The graphs below present impact of fiscal and monetary shocks in the version of the EAGLE model: One with four regions, and one with 15 regions. 

\section{Fiscal Shock}
\begin{figure}[H]
    \centering
    \begin{subfigure}{0.49\textwidth}
        \centering
        \includegraphics[width=\textwidth]{D:/2024-01_eagle/docs/fiscalContributions_Big1.png}
        \caption{Fiscal shock -  4 regions}
        \label{fig:picture1}
    \end{subfigure}
    \hfill
    \begin{subfigure}{0.49\textwidth}
        \centering
        \includegraphics[width=\textwidth]{D:/2024-01_eagle/docs/fiscalContributions_Big12.png}
        \caption{Fiscal shock - 15 regions}
        \label{fig:picture2}
    \end{subfigure}
    \caption{Fiscal shock - Output and inflation, difference from baseline}
    \label{fig:comparison}
\end{figure}

\section{Monetary Shock}
\begin{figure}[H]
    \centering
    \begin{subfigure}{0.45\textwidth}
        \centering
       \includegraphics[width=\textwidth]{D:/2024-01_eagle/docs/monetaryContributions_Big1.png}
        \caption{Fiscal shock - 4 regions}
        \label{fig:picture1}
    \end{subfigure}
    \hfill
    \begin{subfigure}{0.45\textwidth}
        \centering
       \includegraphics[width=\textwidth]{D:/2024-01_eagle/docs/monetaryContributions_Big12.png}
        \caption{Fiscal shock - 15 regions}
        \label{fig:picture2}
    \end{subfigure}
    \caption{Comparison of Picture 1 and Picture 2}
    \label{fig:comparison}
\end{figure}

\end{document}